\documentclass[11pt,a4paper]{article}

\usepackage{amsmath,amssymb}
\usepackage{booktabs}
\usepackage{listings}
\usepackage{xcolor}
\usepackage{hyperref}
\usepackage{geometry}
\usepackage{subfiles}
\geometry{margin=2.5cm}

\definecolor{codebg}{rgb}{0.97,0.97,0.97}
\definecolor{codeframe}{rgb}{0.80,0.80,0.80}

\lstset{
    language=Python,
    basicstyle=\ttfamily\small,
    backgroundcolor=\color{codebg},
    frame=single,
    rulecolor=\color{codeframe},
    columns=fullflexible,
    keepspaces=true,
    showstringspaces=false,
    xleftmargin=1em,
    xrightmargin=1em,
}

\title{\textbf{pysib} --- System Identification Toolbox\\[0.5em]
       \large User Manual\\[0.5em]
       \normalsize Version 0.2.3}
\author{Diego Eckhard\\
        Universidade Federal do Rio Grande do Sul\\
        \texttt{diego@eckhard.com.br}}
\date{}

\begin{document}
\maketitle
\tableofcontents
\newpage

% =========================================================
\section{Overview}
% =========================================================

\textbf{pysib} is a Python toolbox for parameter identification of
discrete-time, single-input single-output (SISO) dynamic systems.
It implements the classical prediction-error estimators for polynomial
input--output model structures: ARX, ARMAX, OE, and Box--Jenkins,
together with auxiliary methods (Stieglitz--McBride, instrumental
variables, correlation error minimization) and filtered variants
designed to improve convergence in the presence of local minima.

\subsection{General model structure}

All model structures supported by pysib are special cases of

\begin{equation}
  y(t) =
  \frac{B(q^{-1})}{A(q^{-1})F(q^{-1})}\,u(t)
  + \frac{C(q^{-1})}{A(q^{-1})D(q^{-1})}\,e(t),
  \label{eq:general}
\end{equation}

where $u(t)$ is the input, $y(t)$ is the output, $e(t)$ is white
noise, and $q^{-1}$ is the backward-shift operator.
The polynomials are written in non-negative powers of $q^{-1}$:
\begin{align*}
  A(q^{-1}) &= 1 + a_1 q^{-1} + \cdots + a_{n_a} q^{-n_a}, \\
  B(q^{-1}) &= b_0 + b_1 q^{-1} + \cdots + b_{n_b-1} q^{-(n_b-1)}, \\
  C(q^{-1}) &= 1 + c_1 q^{-1} + \cdots + c_{n_c} q^{-n_c}, \\
  D(q^{-1}) &= 1 + d_1 q^{-1} + \cdots + d_{n_d} q^{-n_d}, \\
  F(q^{-1}) &= 1 + f_1 q^{-1} + \cdots + f_{n_f} q^{-n_f}.
\end{align*}

Each model structure fixes some polynomials to unity:

\begin{center}
\begin{tabular}{llllll}
  \toprule
  Structure & $A$ & $B$ & $C$ & $D$ & $F$ \\
  \midrule
  ARX    & free & free & $1$  & $1$  & $1$  \\
  ARMAX  & free & free & free & $1$  & $1$  \\
  OE     & $1$  & free & $1$  & $1$  & free \\
  BJ     & $1$  & free & free & free & free \\
  \bottomrule
\end{tabular}
\end{center}

An input delay of $n_k$ samples is incorporated by prepending $n_k$
zeros to the $B$ polynomial.

\subsection{Return convention}

Every identification routine returns a pair \texttt{(theta, m)}:

\begin{itemize}
  \item \texttt{theta}: 1-D \texttt{ndarray} of estimated free parameters,
        in the order listed in the API reference for each function.
  \item \texttt{m}: \texttt{dict} with keys \texttt{A}, \texttt{B},
        \texttt{C}, \texttt{D}, \texttt{F}, each a 1-D \texttt{ndarray}
        of polynomial coefficients with the leading $1$ included where
        applicable.
\end{itemize}

The \texttt{m} dictionary can be passed directly to \texttt{predict}
and \texttt{simulate}, regardless of which estimator produced it.

% =========================================================
\newpage
\section{Installation}
% =========================================================

\subsection{Requirements}

\begin{center}
\begin{tabular}{ll}
  \toprule
  Dependency & Minimum version \\
  \midrule
  Python     & 3.9  \\
  NumPy      & 1.20 \\
  SciPy      & 1.7  \\
  matplotlib & 3.4  \\
  \midrule
  C compiler & GCC $\geq$ 9 or Clang $\geq$ 11 \\
  LAPACK     & Accelerate (macOS) or \texttt{liblapack-dev} (Linux) \\
  \bottomrule
\end{tabular}
\end{center}

The C compiler and LAPACK are required only when building from source,
for example on platforms for which a binary wheel is not available.
On macOS, the Accelerate framework is available by default.
On Linux (Debian/Ubuntu), install the build dependencies with:

\begin{lstlisting}
sudo apt-get install gcc liblapack-dev
\end{lstlisting}

\subsection{Installing via pip}

\begin{lstlisting}
pip install pysib
\end{lstlisting}

On supported platforms, pip installs a binary wheel containing the
compiled extension modules.  If no compatible wheel is available, pip
falls back to building from the source distribution; in that case a
working C compiler and LAPACK must be available before running this
command.

\subsection{Verifying the installation}

\begin{lstlisting}
python -c "import pysib; print('pysib OK')"
\end{lstlisting}

To run the test suite:

\begin{lstlisting}
pip install pytest
python -m pytest tests/ -v
\end{lstlisting}

\subsection{Running examples}

Example scripts are available in the repository under \texttt{examples/}.
They are self-contained and can be run directly from the repository root,
for example:

\begin{lstlisting}
python examples/example_arx.py
\end{lstlisting}

The main user-facing documentation is available at
\url{https://pysib.net/}.

% =========================================================
\newpage
\section{API Reference}
% =========================================================

All functions are accessible directly from the \texttt{pysib} namespace
after \texttt{import pysib}.

% ---------------------------------------------------------
\subsection{Methods without optimization}
% ---------------------------------------------------------

These estimators do not require nonlinear optimization.
ARX and correlation solve a single linear least-squares problem.
IV uses the same structure with an instrument matrix.
SM solves a sequence of linear problems iteratively (100~steps).
They are generally fast, but can be biased when the noise model
does not match the assumed structure.

\subfile{api/arx}
\subfile{api/sm}
\subfile{api/correlation}
\subfile{api/iv}

% ---------------------------------------------------------
\newpage
\subsection{Methods with optimization}
% ---------------------------------------------------------

These estimators minimize the prediction-error criterion
\[
  V(\theta) = \frac{1}{N}\sum_t \varepsilon(t;\theta)^2.
\]
The criterion is nonlinear in $\theta$ because the denominator
polynomials ($F$, or $C$ and $D$) appear in the predictor.
An ARX estimate provides the initial guess in all cases; the optimizer
then runs two phases implemented in a C extension.

\paragraph{Sensitivity signals.}
At each $\theta$, the sensitivity vector
$\psi(t;\theta) = -\partial\varepsilon(t;\theta)/\partial\theta$
is computed by filtering.  For the OE structure, for example,
\[
  \psi^B_j(t) = \frac{u(t-n_k)}{F(q^{-1})} \cdot q^{-j}, \qquad
  \psi^F_j(t) = \frac{\hat{y}(t)}{F(q^{-1})} \cdot q^{-j},
\]
where $\hat{y}(t) = (B/F)\,u(t-n_k)$ is the current simulated
output.  ARMAX and BJ use analogous filtered signals for their
respective polynomials.

\paragraph{Gradient and Gauss--Newton Hessian.}
The gradient and an approximate Hessian are assembled from the
sensitivity signals and the current prediction error $\varepsilon(t)$:
\[
  g_j = -\sum_t \varepsilon(t)\,\psi_j(t), \qquad
  H_{jk} = \sum_t \psi_j(t)\,\psi_k(t).
\]
$H$ is the Gauss--Newton approximation to the Hessian of $V$: it
retains only the outer-product term $\Psi^\top\Psi$ and drops the
second-derivative term (which vanishes near the optimum).

\paragraph{Phase~1 --- steepest descent with smoothed direction.}
The gradient direction is smoothed by an exponential moving average before each step:
\[
  d \leftarrow \tfrac{4}{5}\,d + \tfrac{1}{5}\,g.
\]
The update is $\theta \leftarrow \theta - \alpha\,d/\|d\|$, where
$\alpha$ is adapted per step (multiplied by $1.01$ on a successful
reduction, by $0.99$ otherwise).  This phase runs for up to
$100 \times 100$ gradient evaluations.

\paragraph{Phase~2 --- Gauss--Newton with backtracking.}
The Newton direction $d$ is obtained by solving $H\,d = g$ via
LAPACK \texttt{dposv} (Cholesky factorisation for symmetric positive
definite matrices).  The step length grows linearly with iteration
$i$: $\theta \leftarrow \theta - (i/1000)\,d$.  If the step
increases $V$, the candidate is bisected up to 10 times.  This
phase runs for up to 1\,000 iterations.

\subfile{api/oe}
\subfile{api/armax}
\subfile{api/bj}

% ---------------------------------------------------------
\newpage
\subsection{Filtered estimators}
% ---------------------------------------------------------

The filtered estimators address the sensitivity of the nonlinear
prediction-error criterion to local minima by a cost-function shaping
strategy: both \texttt{u} and \texttt{y} are low-pass filtered before
the optimization, which removes high-frequency content and reshapes the
objective to be more convex in the low-frequency region.  The
optimization is repeated 9 times with progressively wider passbands,
then once more on the original unfiltered data to obtain the final
estimate.

\paragraph{Filter sequence for the filtered estimators.}
A first-order DC-normalized IIR filter is used:
\[
  L_i(q^{-1}) = \frac{1 - a_i}{1 - a_i\,q^{-1}}, \quad i = 1,\ldots,9.
\]
The poles $a_i$ are spaced so that the impulse response decays to
$5\%$ in $\tau_i$ samples, with $\tau_i$ decreasing from 36 to 4:
\[
  a_i = 0.05^{1/\tau_i}, \quad
  \tau_i \in \{36,\,32,\,28,\,24,\,20,\,16,\,12,\,8,\,4\}.
\]
The resulting poles decrease from $a_1 \approx 0.920$ to
$a_9 \approx 0.473$, progressively widening the passband.
The same filter sequence is used by \texttt{oe\_filtered},
\texttt{armax\_filtered}, and \texttt{bj\_filtered}.  The initial
guess for stage~1 is an ARX estimate computed on the most-narrowly
filtered data; each subsequent stage warm-starts from the previous
estimate.

\subfile{api/oe_filtered}
\subfile{api/armax_filtered}
\subfile{api/bj_filtered}

% ---------------------------------------------------------
\newpage
\subsection{Extra functions}
% ---------------------------------------------------------

After identification, three utility functions support model validation
and analysis.
\texttt{predict} and \texttt{simulate} both take the model dictionary
\texttt{m} returned by any estimator.
\texttt{predict} computes the one-step-ahead prediction
$\hat{y}(t\mid t-1)$ using both the input and the observed output;
residuals $\varepsilon(t) = y(t) - \hat{y}(t\mid t-1)$ should be
white and uncorrelated with the input for a correct model.
\texttt{simulate} computes the open-loop output $G(q^{-1})\,u(t)$
using only the input signal; it does not use output observations or
the noise model ($C$, $D$ are ignored).
Both functions can be applied to validation data; they assess
complementary aspects of model quality.
\texttt{plota} is a Monte Carlo diagnostic: given a matrix of
parameter estimates from repeated experiments, it sorts the runs and
plots all parameters together so that the spread and bias of the
estimates can be assessed visually.

\subfile{api/predict}
\subfile{api/simulate}
\subfile{api/plota}

% =========================================================
\newpage
\section{Examples}
% =========================================================

\subsection{ARX}

Based on \texttt{examples/example\_arx.py}.

\begin{lstlisting}
import numpy as np
import matplotlib.pyplot as plt
from scipy.signal import lfilter
import pysib

N = 1000
t = np.arange(N)
u = np.sign(np.sin(2 * np.pi * t / 100))  # square wave

# True model:  y(t) = 0.9*y(t-1) + u(t-1) + e(t)
y = lfilter([0, 1], [1, -0.9], u) + lfilter([1], [1, -0.9], np.random.randn(N))

theta, m = pysib.arx(u, y, na=1, nb=1, nz=1)
print("theta:", theta)

yp = pysib.predict(u, y, m)
ys = pysib.simulate(u, m)

plt.plot(t, y, label="Data")
plt.plot(t, yp, label="Prediction")
plt.plot(t, ys, label="Simulation")
plt.legend(); plt.show()
\end{lstlisting}

\subsection{OE and Stieglitz--McBride}

Based on \texttt{examples/example\_oe.py} and \texttt{examples/example\_sm.py}.

\begin{lstlisting}
import numpy as np
import matplotlib.pyplot as plt
from scipy.signal import lfilter
import pysib

N = 100
t = np.arange(N)
u = np.sign(np.sin(2 * np.pi * t / 100))

# True OE model:  y(t) = u(t-1) / (1 - 0.9*z^-1) + e(t)
y = lfilter([0, 1], [1, -0.9], u) + 0.01 * np.random.randn(N)

theta_oe, m_oe = pysib.oe(u, y, nb=1, nf=1, nz=1)
theta_sm, m_sm = pysib.sm(u, y, nb=1, nf=1, nz=1)
print("OE theta:", theta_oe)
print("SM theta:", theta_sm)

yp = pysib.predict(u, y, m_oe)
ys = pysib.simulate(u, m_oe)

plt.plot(t, y, label="Data")
plt.plot(t, yp, label="Prediction (OE)")
plt.plot(t, ys, label="Simulation (OE)")
plt.legend(); plt.show()
\end{lstlisting}

\subsection{ARMAX}

Based on \texttt{examples/example\_armax.py}.

\begin{lstlisting}
import numpy as np
import matplotlib.pyplot as plt
from scipy.signal import lfilter
import pysib

N = 1000
t = np.arange(N)
u = np.sign(np.sin(2 * np.pi * t / 100))

# True ARMAX:  y = u(t-1)/(1-0.9z^-1) + (1-0.5z^-1)/(1-0.9z^-1)*e(t)
y = (lfilter([0, 1], [1, -0.9], u)
   + lfilter([1, -0.5], [1, -0.9], np.random.randn(N)))

theta, m = pysib.armax(u, y, na=1, nb=1, nc=1, nz=1)
print("theta:", theta)

yp = pysib.predict(u, y, m)
ys = pysib.simulate(u, m)

plt.plot(t, y, label="Data")
plt.plot(t, yp, label="Prediction")
plt.plot(t, ys, label="Simulation")
plt.legend(); plt.show()
\end{lstlisting}

\subsection{Box--Jenkins}

Based on \texttt{examples/example\_bj.py}.

\begin{lstlisting}
import numpy as np
import matplotlib.pyplot as plt
from scipy.signal import lfilter
import pysib

N = 1000
t = np.arange(N)
u = np.sign(np.sin(2 * np.pi * t / 100))

# True BJ:  y = u(t-1)/(1-0.9z^-1) + (1-0.5z^-1)/(1-1.5z^-1+0.7z^-2)*e(t)
y = (lfilter([0, 1], [1, -0.9], u)
   + lfilter([1, -0.5], [1, -1.5, 0.7], 0.05 * np.random.randn(N)))

theta, m = pysib.bj(u, y, nb=1, nc=1, nd=2, nf=1, nz=1)
print("theta:", theta)

yp = pysib.predict(u, y, m)
ys = pysib.simulate(u, m)

plt.plot(t, y, label="Data")
plt.plot(t, yp, label="Prediction")
plt.plot(t, ys, label="Simulation")
plt.legend(); plt.show()
\end{lstlisting}

\subsection{Correlation error minimization}

Based on \texttt{examples/example\_correlation.py}.

\begin{lstlisting}
import numpy as np
import matplotlib.pyplot as plt
from scipy.signal import lfilter
import pysib

N = 1000
t = np.arange(N)
u = np.sign(np.sin(2 * np.pi * t / 100))

# True ARX: y(t) = 0.9*y(t-1) + u(t-1) + e(t)
y = lfilter([0, 1], [1, -0.9], u) + np.random.randn(N)

theta_ls,   m_ls   = pysib.arx(u, y, na=1, nb=1, nz=1)
theta_corr, m_corr = pysib.correlation(u, y, na=1, nb=1, nz=1, M=30)
print("ARX        :", theta_ls)
print("Correlation:", theta_corr)

yp = pysib.predict(u, y, m_corr)
ys = pysib.simulate(u, m_corr)

plt.plot(t, y, label="Data")
plt.plot(t, yp, label="Prediction (correlation)")
plt.plot(t, ys, label="Simulation (correlation)")
plt.legend(); plt.show()
\end{lstlisting}

\subsection{Instrumental variables}

Based on \texttt{examples/example\_iv.py}.

\begin{lstlisting}
import numpy as np
from scipy.signal import lfilter
import pysib

N = 1000
t = np.arange(N)
u = np.sign(np.sin(2 * np.pi * t / 100))

# Two independent noise realizations, same input
y1 = lfilter([0, 1], [1, -0.9], u) + np.random.randn(N)
y2 = lfilter([0, 1], [1, -0.9], u) + np.random.randn(N)

theta_arx, _ = pysib.arx(u, y1, na=1, nb=1, nz=1)
theta_iv,  _ = pysib.iv(u, y1, y2, na=1, nb=1, nz=1)
print("ARX:", theta_arx)
print("IV :", theta_iv)
\end{lstlisting}

\subsection{OE with filtered initialization}

Based on \texttt{examples/example\_oe\_filtered.py}.

\begin{lstlisting}
import numpy as np
import matplotlib.pyplot as plt
from scipy.signal import lfilter
import pysib

N = 100
t = np.arange(N)
u = np.sign(np.sin(2 * np.pi * t / 100))

y = lfilter([0, 1], [1, -0.9], u) + 0.1 * np.random.randn(N)

theta_oe,   m_oe   = pysib.oe(u, y, nb=1, nf=1, nz=1)
theta_filt, m_filt = pysib.oe_filtered(u, y, nb=1, nf=1, nz=1)
print("OE         :", theta_oe)
print("OE filtered:", theta_filt)

yp = pysib.predict(u, y, m_filt)
ys = pysib.simulate(u, m_filt)

plt.plot(t, y, label="Data")
plt.plot(t, yp, label="Prediction (filtered)")
plt.plot(t, ys, label="Simulation (filtered)")
plt.legend(); plt.show()
\end{lstlisting}

\subsection{ARMAX with filtered initialization}

Based on \texttt{examples/example\_armax\_filtered.py}.

\begin{lstlisting}
import numpy as np
from scipy.signal import lfilter
import pysib

N = 1000
t = np.arange(N)
u = np.sign(np.sin(2 * np.pi * t / 100))

y = (lfilter([0, 1], [1, -0.9], u)
   + lfilter([1, -0.5], [1, -0.9], np.random.randn(N)))

theta_armax, _ = pysib.armax(u, y, na=1, nb=1, nc=1, nz=1)
theta_filt,  _ = pysib.armax_filtered(u, y, na=1, nb=1, nc=1, nz=1)
print("ARMAX         :", theta_armax)
print("ARMAX filtered:", theta_filt)
\end{lstlisting}

\subsection{Box--Jenkins with filtered initialization}

Based on \texttt{examples/example\_bj\_filtered.py}.

\begin{lstlisting}
import numpy as np
from scipy.signal import lfilter
import pysib

N = 1000
t = np.arange(N)
u = np.sign(np.sin(2 * np.pi * t / 100))

y = (lfilter([0, 1], [1, -0.9], u)
   + lfilter([1, -0.5], [1, -1.5, 0.7], 0.05 * np.random.randn(N)))

theta_bj,   _ = pysib.bj(u, y, nb=1, nc=1, nd=2, nf=1, nz=1)
theta_filt, _ = pysib.bj_filtered(u, y, nb=1, nc=1, nd=2, nf=1, nz=1)
print("BJ         :", theta_bj)
print("BJ filtered:", theta_filt)
\end{lstlisting}

\subsection{Monte Carlo plot}

Based on \texttt{examples/example\_monte\_carlo.py}.

\begin{lstlisting}
import numpy as np
import matplotlib.pyplot as plt
from scipy.signal import lfilter
import pysib

M = 100
N = 100
T = np.full((2, M), np.nan)

for run in range(M):
    t = np.arange(N)
    u = np.sign(np.sin(2 * np.pi * t / 100))
    y = lfilter([0, 1], [1, -0.9], u) + 10 * np.random.randn(N)
    theta, _ = pysib.oe(u, y, nb=1, nf=1, nz=1)
    T[:, run] = theta

T_sorted = pysib.plota(T, a=0)
plt.title("OE Monte-Carlo (100 runs)")
plt.show()
\end{lstlisting}

% =========================================================
\newpage
\section{Parameter selection guide}
% =========================================================

\subsection{Model orders}

\begin{description}
  \item[\texttt{na}] Number of $A$ parameters ($a_1,\ldots,a_{n_a}$).
        Start with \texttt{na = 1} or \texttt{2} for first- and
        second-order systems.

  \item[\texttt{nb}] Number of $B$ parameters ($b_0, \ldots, b_{n_b-1}$).
        \texttt{nb = 1} gives a single gain $b_0$; the degree of $B$
        is $n_b - 1$.

  \item[\texttt{nf}] Number of $F$ parameters ($f_1,\ldots,f_{n_f}$),
        equal to the degree of $F$ (poles of the noise-free model in
        OE and BJ).  Typically equal to the system order.

  \item[\texttt{nc}] Number of $C$ parameters ($c_1,\ldots,c_{n_c}$),
        equal to the degree of $C$ (noise MA polynomial in ARMAX and
        BJ).  Start with \texttt{nc = 1} and increase only if
        residuals remain correlated.

  \item[\texttt{nd}] Number of $D$ parameters ($d_1,\ldots,d_{n_d}$),
        equal to the degree of $D$ (noise AR polynomial in BJ only).
        Start with \texttt{nd = 1}.

  \item[\texttt{nz}] Input delay in samples.  \texttt{nz = 1} means the
        input affects the output one step later ($b_0$ multiplies
        $u(t-1)$).  Set \texttt{nz} equal to the known or estimated
        dead time of the process.
\end{description}

\subsection{Which estimator to use}

The central question is \emph{where the noise enters the system}.
All model structures are special cases of
$y = G\,u + H\,e$ where $G = B/(AF)$ is the plant and
$H = C/(AD)$ is the noise model.  Different assumptions about $H$
lead to different estimators.

\paragraph{Process noise (ARX, ARMAX).}
If the disturbance enters at the plant input, it is filtered by the
plant before reaching the output.  The output noise therefore has
the same poles as the plant, giving the ARX noise model
$H = 1/A$.  In this case ordinary least squares on the ARX
structure is consistent.  If the residuals from ARX are correlated
(the noise has additional moving-average dynamics), extend to ARMAX
by adding the $C$ polynomial, making $H = C/A$.

\paragraph{Measurement noise (OE, BJ).}
If the disturbance is additive at the output (sensor or measurement
noise), the noise is independent of the plant dynamics.  The OE
structure captures this with $H = 1$ (white noise).  If the
measurement noise is itself colored, BJ provides the most general
description: $H = C/D$ with poles and zeros independent of the
plant $G = B/F$.  Applying ARX to data with additive output noise
yields biased estimates because the lagged output regressors are
then correlated with the noise.

\paragraph{Moment-based methods for the ARX structure (IV, correlation).}
When ARX ordinary least squares is expected to be biased, two
moment-based alternatives avoid nonlinear optimization entirely.
The \emph{instrumental variable} method requires two independent
experiments with the same input: it uses the output of the second
experiment as an instrument for the first, yielding a consistent
estimate because the instrument shares the input-driven component
but has independent noise.
The \emph{correlation} method works with a single experiment: it
finds $\theta$ such that the prediction error is uncorrelated with
the instrument $z$ at $M$ lags.  By default $z = u$, which is a
valid instrument whenever the input is uncorrelated with the noise
(open-loop operation).  An external instrument can be supplied when
this condition does not hold.  Using $M > n_a + n_b$ lags
over-determines the system and improves robustness.

\vspace{0.5em}
\begin{center}
\begin{tabular}{p{0.38\textwidth}p{0.13\textwidth}p{0.35\textwidth}}
  \toprule
  Situation & Estimator & Notes \\
  \midrule
  Quick estimate; disturbance enters at plant input (process noise) &
  ARX &
  Closed-form; consistent when $H = 1/A$.
  Biased under additive output noise. \\
  \addlinespace
  Additive output noise (measurement or sensor noise) &
  OE &
  Consistent plant model $B/F$; nonlinear optimization. \\
  \addlinespace
  OE structure; faster computation preferred &
  SM &
  Iterative ARX steps (100 iterations); cheaper but may
  converge to a local minimum. \\
  \addlinespace
  ARX residuals are correlated; noise shares plant denominator &
  ARMAX &
  Extends ARX with MA noise $C/A$; nonlinear optimization. \\
  \addlinespace
  Plant and noise have independent pole structures &
  BJ &
  Fully separate plant $B/F$ and noise $C/D$;
  most parameters; needs most data. \\
  \addlinespace
  Two independent experiments with the same input available &
  IV &
  Consistent under ARX noise structure;
  no nonlinear optimization. \\
  \addlinespace
  Single experiment; input uncorrelated with noise (open-loop);
  or external instrument available &
  Correlation &
  Consistent ARX estimate without a second experiment;
  $M > n_a + n_b$ lags improve robustness. \\
  \bottomrule
\end{tabular}
\end{center}

\subsection{Filtered vs.\ standard estimators}

The prediction-error criterion $V(\theta)$ is nonlinear in $\theta$
whenever denominator polynomials are free ($F$ in OE, $C$ or $D$
in ARMAX and BJ).  Nonlinear optimizers — including the
steepest-descent and Gauss--Newton algorithm implemented here —
can converge to a local minimum rather than the global one,
particularly when the system has slow dynamics or the data record
is short.

The filtered variants reduce the likelihood of this happening by
reshaping the cost function before the final optimization.
Low-pass filtering the data at progressively wider bandwidths
produces a sequence of modified criteria, each of which tends to
have fewer local minima than the full-bandwidth criterion.
Each stage is warm-started from the solution of the previous one,
so the optimizer is guided toward a region close to the global
minimum before the final unfiltered step.
This does not \emph{guarantee} convergence to the global minimum,
but it substantially reduces the probability of failure in
practice.

The computational cost of a filtered variant is approximately
$10\times$ that of the corresponding standard estimator
(9 filtered stages plus the final unfiltered run, each requiring
a full optimization).  Use the standard estimator first; switch
to the filtered version when the standard one returns unstable
polynomials, produces estimates that vary across runs, or fails
to reduce $V$ to a plausible level.

\end{document}
